\newcommand{\tipoRequerimento}{cobrança indevida sobre enquadramento tarifário incorreto}
\section{DOS FATOS E DOS DIREITOS}

\begin{enumerate}[resume]
\item
Verificou-se que a UC \unidadeConsumidora{} está cadastrada com enquadramento tarifário errado, gerando, desta forma, prejuízos ao consumidor, por erro de cadastro por parte da concessionária.
\item  
Assim, a UC \unidadeConsumidora{} deve ser cadastrada conforme art. 189, incisos I e II, da Resolução Normativa 1000/2021 da Aneel, vale dizer, deve ser enquadrada na classe iluminação pública (B4A).  
\item 
A seguir transcreve-se o art. 189 da citada resolução:

\Citacao{Art. 189. Deve ser classificada na classe iluminação pública a unidade consumidora destinada exclusivamente à prestação do serviço público de iluminação pública, de responsabilidade do poder público municipal ou distrital ou daquele que receba essa delegação, com o objetivo de iluminar:   

I - vias públicas destinadas ao trânsito de pessoas ou veículos, tais como ruas, avenidas, logradouros, caminhos, passagens, passarelas, túneis, estradas e rodovias; e   

II - bens públicos destinados ao uso comum do povo, tais como abrigos de usuários de transportes coletivos, praças, parques e jardins, ainda que o uso esteja sujeito a condições estabelecidas pela administração, inclusive o cercamento, a restrição de horários e a cobrança.  }

\item 
A seguir transcreve-se o art. 53-O da REN 414/2010 (revogada pela REN 1.000/2021) da Aneel:

\Citacao{Art. 53. Na classe iluminação pública enquadram-se as unidades consumidoras destinadas exclusivamente para a prestação do serviço público de iluminação pública, de responsabilidade do Público Municipal ou Distrital, ou ainda daquele que receba essa delegação, com o objetivo de iluminar: (Incluído pela REN ANEEL 800, de 19.12.2017).   

I - vias públicas destinadas ao trânsito de pessoas ou veículos, tais como ruas, avenidas, logradouros, caminhos, passagens, passarelas, túneis, estradas e rodovias; e (Incluído pela REN ANEEL 800, de 19.12.2017)   

II - bens públicos destinados ao uso comum do povo, tais como abrigos de usuários de transportes coletivos, praças, parques e jardins, ainda que o uso esteja sujeito a condições estabelecidas pela administração, inclusive o cercamento, a restrição de horários e a cobrança. (Incluído pela REN ANEEL 800, de 19.12.2017). (Grifou-se) }

\item
Logo, sempre que a iluminação tiver por finalidade iluminar bens de uso comum do povo deverá ser classificada na classe iluminação pública.   
\item 
O Código Civil de 2002 divide os bens públicos, segundo à sua destinação, em três categorias: bens de uso comum do povo ou de domínio público, bens de uso especial ou do patrimônio administrativo indisponível e bens dominicais ou do patrimônio disponível.   
\item
Os bens de uso comum do povo são os bens que se destinam à utilização geral pela coletividade (como por exemplo, ruas, praças, estradas, quadras de esporte, campos de futebol, cemitérios, jardins, avenidas, ginásios etc). Vale dizer, são todos aqueles bens de utilização concorrente de toda a comunidade, usados livremente pela população, o que não significa “de graça” e sim, que não dependem de prévia autorização do Poder Público para sua utilização.   
\item
Vale destacar que a relação do art. 99 do Código Civil, é apenas exemplificativa:

\Citacao{Art. 99. São bens públicos:   

I - os de uso comum do povo, tais como rios, mares, estradas, ruas e praças;

II - os de uso especial, tais como edifícios ou terrenos destinados a serviço ou estabelecimento da administração federal, estadual, territorial ou municipal, inclusive os de suas autarquias;

III - os dominicais, que constituem o patrimônio das pessoas jurídicas de direito público, como objeto de direito pessoal, ou real, de cada uma dessas entidades.}

\item 
Podemos demonstrar vários exemplos de jurisprudência que citam quadras de esportes, campos de futebol, dentre outros, como bem público de uso comum do povo. 

\Citacao{TSE - RESPE: 5067820126160073 Pato Branco/PR 271092013, Relator: Min. Maria Thereza Rocha De Assis Moura, Data de Julgamento: 25/02/2015, Data de Publicação: DJE - Diário de justiça eletrônico - 04/03/2015 - Página 135-137. }

\end{enumerate}

\section{DOS PEDIDOS}

\begin{enumerate}[resume]
\item 
Diante do exposto, requer que a Concessionária:

\begin{enumerate}
\item 
Efetue vistoria em tal UC e proceda com a correção do enquadramento tarifário da mesma em conformidade com o estabelecido na Resolução Normativa Nº 1.000/2021 da ANEEL.
\item 
Efetue a devolução de todos os valores cobrados a maior, referente aos últimos 10 (dez) anos. Conforme o Despacho Nº 2.006 da ANEEL de 08 de julho de 2024, observando o prazo prescricional de 10 anos, conforme previsto no art. 205 do Código Civil.
\item 
Devolva os valores cobrados a maior em dobro, atualizados monetariamente pelo IPCA e juros de mora de 1\% ao mês calculados pro \textit{rata die}, conformidade com o estabelecido na Resolução Normativa Nº 1.000/2021 da ANEEL.   
\item 
Que apresente as seguintes informações: Data de ligação da UC, se houve modificação de tarifa entre a data da ligação e a data da resposta.   
\item 
Que em caso de indeferimento apresente os argumentos e as provas (comprovações) para não enquadramento (fotos, carga instalada etc).   
\item 
Proceda a análise e atendimento a este requerimento no prazo máximo de 10 (dez) dias úteis, para informar o resultado da classificação tarifária no caso de necessidade de vistoria, em conformidade com o estabelecido na Resolução Normativa Nº 1.000/2021 da ANEEL.  
\item 
Proceda com a devolução do valor integral, por meio de depósito bancário, na conta corrente do Município, conforme estabelece o § 7º, do art. 323 da REN Nº 1.000/2021 da ANEEL, comunicando por escrito o resultado do deferimento desta reclamação, por meio de e-mail.

\end{enumerate}
\end{enumerate}